\section{Kiểm tra bổ sung và giới hạn triển khai}
\label{muc:kiemtraphu}

Các phép kiểm trong mục này trả lời ba phản biện đối với kết quả chính: giảm dịch chuyển
có phụ thuộc cách đo hay không, các cơ chế bảo vệ có thực sự khác nhau trên bộ sinh đang
dùng hay không, và chi phí phục vụ có phải do giới hạn tìm kiếm. Đây là các phân tích bổ
sung trên những lần chạy đã có. Chúng không thay thế phép so sánh chính và không mở rộng
phạm vi kết luận sang hệ điều phối khác.

\subsection{Kết quả phụ thuộc đại lượng được đo}


Phân tích ban đầu dùng chín đại lượng mô tả: tổng tuyến tính, tổng bậc hai, hai tổng có
trọng số theo thời gian còn lại, số lần dời, cực đại và hai phân vị của gánh nặng mỗi khách,
cùng số khách bị dời. Bảng~\ref{bang:kiemtra-thuocdo} giữ bốn đại lượng có ý nghĩa trực
tiếp. Tổng và cực đại trong bảng cộng hai chiều giờ đón và giờ trả; chúng không phải riêng
sổ cái giờ trả được chặn ở trần 30 giây. ``Một lần dời'' là một lần công bố có tổng phần
dịch chuyển do quyết định lớn hơn không, còn ``khách bị dời'' đếm định danh khách riêng biệt.

\begin{table}[htbp]
\centering
\caption{Bốn đại lượng mô tả dịch chuyển do quyết định trên cùng tập 20 ô.
Hai cột thời gian có đơn vị giây; ngân sách trong tên cấu hình chỉ áp cho giờ trả.}
\label{bang:kiemtra-thuocdo}
\small
\begin{tabular}{@{}lrrrr@{}}
\toprule
\textbf{Cấu hình} & \textbf{Tổng} & \textbf{Cực đại/khách} & \textbf{Lần dời} & \textbf{Khách bị dời}\\
\midrule
B1 & 74.443 & 840{,}9 & 479 & 403\\
Không trần giờ trả & 22.485 & 595{,}3 & 155 & 142\\
Trần 300 s & 15.809 & 290{,}8 & 135 & 128\\
Trần 120 s & 5.063 & 119{,}0 & 79 & 76\\
Trần 60 s & 1.109 & 59{,}0 & 34 & 34\\
C1, trần 30 s & 128 & 27{,}6 & 11 & 11\\
\bottomrule
\end{tabular}
\end{table}

Bốn đại lượng đều giảm theo thứ tự cấu hình được trình bày; các đại lượng bổ sung trong
phân tích chín thước đo cho cùng hướng trên tập ô này. Tuy nhiên, độ lớn của thay đổi phụ
thuộc câu hỏi: tổng dịch chuyển, số lần dời và mức một khách chịu không thể dùng thay cho
nhau. Kết quả này chỉ mô tả thứ hạng tổng hợp trên sáu cấu hình, không chứng minh thứ hạng
ở mọi ô, mọi định nghĩa bất ổn hay mọi dữ liệu. Đặc biệt, tính đơn điệu của các tổng dịch
chuyển không kéo theo tính đơn điệu của số khách phục vụ.

Một khác biệt khác là đo theo hai đầu hay theo cả chuỗi. Phép đối chiếu lấy độ muộn dương
của giờ trả công bố cuối so với giờ trả công bố đầu, và so với tổng độ lớn các thay đổi
giữa các lần công bố. Bảng~\ref{bang:kiemtra-hai-dau} cho kết quả trên hai nhánh gốc.
Đại lượng hai đầu liên quan đến thước đo của Fielbaum và Alonso-Mora
\cite{fielbaum2020}, nhưng ở đây dùng giờ công bố cuối làm mốc cuối. Nó không thay thế
việc nối bản ghi với giờ thực hiện khi cần đo độ tin cậy thực tế.

\begin{table}[htbp]
\centering
\caption{Độ lệch hai đầu và độ dài chuỗi giờ trả trên cùng dữ liệu, tính bằng giây.
Hai đại lượng đều tính theo giờ khách thấy, khác với tổng phần do quyết định.}
\label{bang:kiemtra-hai-dau}
\small
\begin{tabular}{@{}lrr@{}}
\toprule
\textbf{Đại lượng} & \textbf{B1} & \textbf{C1}\\
\midrule
Tổng độ muộn dương, công bố cuối so với đầu & 64.814 & 126\\
Tổng biến động trên chuỗi công bố giờ trả & 71.529 & 6.822\\
\bottomrule
\end{tabular}
\end{table}

Ở B1, hai tổng gần nhau vì phần dịch chuyển lớn do quyết định chủ yếu làm muộn giờ.
Ở C1, biến động trên chuỗi vẫn lớn hơn độ lệch hai đầu nhiều lần. Phần chênh gần như
toàn là thay đổi khi giữ nguyên tuyến, đi theo cả hai chiều và triệt tiêu ở hai đầu;
bộ sinh không tạo những sửa chữa kéo giờ khách cũ sớm lên. Vì thế không được diễn giải
bảng thành bằng chứng rằng thước đo hai đầu bỏ sót một lượng lớn hành vi có hại của
điều phối. Hai cách đo trả lời hai câu hỏi khác nhau; cả hai cần được đọc cùng phép
phân rã vận hành.

\subsection{Bảo vệ giờ đón phụ thuộc khả năng của bộ sinh}


Bảo đảm một chiều cho phép giờ đón tốt lên nhưng cấm xấu đi. Nó là một cơ chế đã có
trong văn liệu, không phải đóng góp mới của khóa luận. Trong bộ sinh chèn thuần đang
dùng, thêm điểm dừng không tạo một phương án kéo giờ đón khách cũ sớm hơn mốc giữ tuyến.
Do đó quyền cải thiện mà bảo đảm một chiều cho phép ít có tác dụng trong thiết lập này.

Phép thử cơ chế này dùng một bộ chạy riêng vì bộ chạy gốc chưa hỗ trợ trường cấu
hình tương ứng. Hai nhánh đối chứng được chạy trên đủ 20 ô và tái lập kết quả của
bộ chạy gốc trước khi đọc nhánh mới. Nhánh bảo đảm một chiều và nhánh khóa giờ đón
đều phục vụ $1.581$ khách, có tổng dịch chuyển do quyết định $128$ giây, $11$ lần
dời và cực đại $27{,}6$ giây trên một khách. Đây là sự tương đương quan sát được
trong bộ sinh và cấu hình đã thử, không phải một phép tái lập đầy đủ thuật toán
của nghiên cứu trước, cũng không phải khẳng định hai cơ chế luôn tương đương.

Hàng rào đóng băng 600 giây cũng cho cùng các kết quả chính với khóa giờ đón.
Nhưng tầm đặt chỗ của kịch bản chỉ là 600 giây: một hàng rào phủ toàn tầm đặt chỗ
không có cơ hội mở. Vì vậy phép thử không đánh giá lợi ích của hàng rào dài trên
yêu cầu đặt trước xa hơn. Các kết quả này cho thấy phải kiểm cả cơ chế bảo vệ và
khả năng bộ sinh tạo phương án sử dụng phần tự do còn lại.

\subsection{Giới hạn tìm kiếm và các phép thử âm}
\label{muc:bosinh}

Trong 40 lần chạy của hai nhánh H6 gốc, bằng chứng sinh phương án không ghi nhận
lần nào cạn ngân sách tìm kiếm, chạm trần ứng viên, bỏ ứng viên khả thi do trần
hay bỏ nhánh chưa khai triển. Các số không này loại được một giải thích cụ thể:
chênh lệch chính không đến từ việc một nhánh bị cắt sớm bởi các giới hạn đã cấu
hình. Chúng không chứng minh tìm kiếm tối ưu trong toàn bộ không gian DARP.
Tập được xét vẫn do bộ sinh chèn, khóa phân công và các kiểm tra vật lý xác định.

Tập vị trí chèn thực sự khá nhỏ. Nếu chỉ lấy các kỷ nguyên có quyết định nhận
hoặc từ chối khách, tỉ lệ tập ứng viên của một xe có ít nhất hai vị trí chèn là
$2{,}48\%$ ở B1 và $2{,}95\%$ ở C1. Phải dùng mẫu số này thay cho mọi kỷ nguyên,
vì các bước chỉ cập nhật tiến độ xe thường chỉ có phương án giữ tuyến. Không
chạm trần ứng viên còn có nghĩa việc đổi thứ tự ưu tiên vị trí chèn không làm
đổi đầu vào bộ giải trong những lần chạy này. Tuy nhiên, điều đó không chứng
minh cắt tỉa sớm không thể giảm thời gian tính ở một triển khai khác.

Bảng~\ref{bang:kiemtra-am} tóm tắt những can thiệp bổ sung có liên quan đến cách
đọc kết quả. Chỉ giữ các trường hợp cho biết giới hạn của mô hình hoặc loại
một giải thích về chi phí phục vụ; không dùng chúng như nhật ký thử nghiệm.

\begin{table}[htbp]
\centering
\caption{Các phép thử bổ sung và giới hạn của điều có thể suy ra.}
\label{bang:kiemtra-am}
\small
\begin{tabular}{@{}>{\raggedright\arraybackslash}p{2.9cm}>{\raggedright\arraybackslash}p{4.1cm}>{\raggedright\arraybackslash}p{7.8cm}@{}}
\toprule
\textbf{Can thiệp} & \textbf{Kết quả quan sát} & \textbf{Điều học được và giới hạn}\\
\midrule
Ưu tiên vị trí chèn & Không có ứng viên bị bỏ do trần trong tập ô gốc &
Đổi thứ tự sinh không đổi tập đưa vào bộ giải ở thiết lập này; không kết luận về
tập ô lớn hoặc có giới hạn tính toán chặt.\\
Ngưỡng cảnh báo 60 s, trần cứng 120 s & Trùng chuỗi lựa chọn với đối chứng ở 20/20 ô &
Tầng cảnh báo không đổi quyết định. Bằng chứng hiện có chưa xác định nguyên nhân
xếp hạng; đây là kết quả âm trong tập ô, không phải chứng minh cơ chế luôn vô hiệu.\\
Sắp xếp lại trên cùng xe & 18/20 ô gặp \texttt{LOCKED\_LEG\_MISMATCH} &
Bộ mô phỏng khóa chặng đang chạy khác với thời điểm đóng băng trong mô hình của
phép thử. Không có phép so hiệu quả hợp lệ cho cả tập ô; không suy ra sắp xếp lại
là bất khả thi về thuật toán.\\
\bottomrule
\end{tabular}
\end{table}

Phép thử ngưỡng cảnh báo còn cho một chẩn đoán triển khai đáng giữ: tầng mức sử
dụng ngân sách tệ nhất nhận cùng giá trị $1.000.000$ ở các quyết định có tầng này,
do cách xử lý chiều có trần cứng bằng không. Tầng đó luôn hòa trong dữ liệu đã
đọc, nên không giải thích được việc tầng cảnh báo không đổi lựa chọn. Những tầng
khác hoặc việc các ứng viên cùng mức vượt cảnh báo vẫn là các giải thích chưa
được phân biệt, vì chứng cứ không ghi đầy đủ điểm mọi ứng viên còn cạnh tranh.

Các phép kiểm bổ sung củng cố cách đọc hẹp của kết quả chính: các cấu hình đã
thực hiện những hành vi được mô tả, và giới hạn tìm kiếm đã ghi không giải thích
chênh lệch H6. Chúng đồng thời cho thấy bảo đảm, bộ sinh, trạng thái chặng đang
chạy và luật phục hồi cần có hợp đồng thống nhất. Một kết quả âm của tổ hợp hiện
tại không đủ để bác bỏ các thiết kế điều phối có khả năng sửa tuyến hoặc tái gán.
